% #1 center of cube
% #2 edge length
\newcommand\getCornersOfCube[2]{%
    \divScalar[halfEdge]{#2}{2}
    \newVec{tmpVec}{-\halfEdge,-\halfEdge,-\halfEdge, 0}
    \addVec[pointA]{#1, \tmpVec}
    \newVec{tmpVec}{-\halfEdge,-\halfEdge,\halfEdge, 0}
    \addVec[pointB]{#1, \tmpVec}
    \newVec{tmpVec}{-\halfEdge,\halfEdge,-\halfEdge, 0}
    \addVec[pointC]{#1, \tmpVec}
    \newVec{tmpVec}{-\halfEdge,\halfEdge,\halfEdge, 0}
    \addVec[pointD]{#1, \tmpVec}
    \newVec{tmpVec}{\halfEdge,-\halfEdge,-\halfEdge, 0}
    \addVec[pointE]{#1, \tmpVec}
    \newVec{tmpVec}{\halfEdge,-\halfEdge,\halfEdge, 0}
    \addVec[pointF]{#1, \tmpVec}
    \newVec{tmpVec}{\halfEdge,\halfEdge,-\halfEdge, 0}
    \addVec[pointG]{#1, \tmpVec}
    \newVec{tmpVec}{\halfEdge,\halfEdge,\halfEdge, 0}
    \addVec[pointH]{#1, \tmpVec}
}

% #1 suffix for points
% #2 first mat, for example: trans mat
% #3 second mat, for example: scale mat
% #4 third mat, for example: rotation mat
\newcommand\getShapeHouse[4][]{%
    \mulMat[fsMat]{#2}{#3}
    \mulMat[fstMat]{\fsMat}{#4}
    \newVec{tmpVec}{-0.5, -0.5, 1}
    \transformVec[pointA#1]{\fstMat}{\tmpVec}
    \newVec{tmpVec}{-0.5, 0.5, 1}
    \transformVec[pointB#1]{\fstMat}{\tmpVec}
    \newVec{tmpVec}{0.5, 0.5, 1}
    \transformVec[pointC#1]{\fstMat}{\tmpVec}
    \newVec{tmpVec}{0.5, -0.5, 1}
    \transformVec[pointD#1]{\fstMat}{\tmpVec}
    \newVec{tmpVec}{0.0, 1.0, 1}
    \transformVec[pointE#1]{\fstMat}{\tmpVec}
}